%% ma: L12-B08-0056
%% lop: 12
%% chuong: 2
%% bai: 8
%% dang: 12.8.33
%% loai: DS
%% muc: [NB, TH, TH, NB]
%% dap_an: "ĐSSĐ"
%% nguon: "LaTeX giáo viên gửi 21/10/2026 (M05, Vectơ và hệ trục tọa độ trong không gian, 107 câu), Phần II Câu 11"
%% kien_thuc: "Bài 8 (tọa độ vectơ, tích vô hướng, trọng tâm tứ diện)"
%% trang_thai: cho-duyet
%% ly_do: "Dạng mới đề xuất 12.8.33"
\begin{ex}
Trong không gian $Oxyz$, cho tứ diện $ABCD$ với $A(2;1;0)$, $B(1;1;3)$, $C(2;-1;3)$, $D(1;-1;0)$.
\choiceTF
{\True Tọa độ vectơ $\overrightarrow{AB}$ là $(-1;0;3)$.}
{Tam giác $ABC$ là tam giác vuông tại $B$.}
{Vectơ $\overrightarrow{AB}$ và vectơ $\overrightarrow{CD}$ vuông góc với nhau.}
{\True Trọng tâm của tứ diện $ABCD$ là điểm $G\left(\dfrac32;0;\dfrac32\right)$.}
\loigiai{a) $\overrightarrow{AB}=(-1;0;3)$. Đúng.

b) $\overrightarrow{BA}=(1;0;-3)$, $\overrightarrow{BC}=(1;-2;0)$, $\overrightarrow{BA}\cdot\overrightarrow{BC}=1\ne0$. Sai.

c) $\overrightarrow{CD}=(-1;0;-3)$, $\overrightarrow{AB}\cdot\overrightarrow{CD}=1-9=-8\ne0$. Sai.

d) $G\left(\dfrac{2+1+2+1}4;\dfrac{1+1-1-1}4;\dfrac{0+3+3+0}4\right)=\left(\dfrac32;0;\dfrac32\right)$. Đúng.}
\end{ex}
